\documentclass[11pt]{article}

% ---------- Packages ----------
\usepackage[a4paper,margin=0.85in]{geometry}
\usepackage{amsmath,amssymb,mathtools}
\usepackage{tikz-cd}
\usepackage{graphicx}
\usepackage{caption}
\usepackage{hyperref}

% ---------- Macros ----------
\newcommand{\Can}{\mathsf{Can}}
\newcommand{\Coeff}{\mathsf{Coeff}}
\newcommand{\CtwoS}{\mathsf{C2S}}
\newcommand{\StwoC}{\mathsf{S2C}}
\newcommand{\EvalMod}{\mathsf{EvalMod}}
\newcommand{\Split}{\mathsf{Split}}
\newcommand{\Merge}{\mathsf{Merge}}
\newcommand{\BTS}{\mathsf{BTS}}

\begin{document}

\title{Core Commutative Diagrams for Coefficient-Split CKKS Bootstrapping}
\author{}
\date{}
\maketitle

This note contains two commutative diagrams for the ideal normalized
coefficient-split CKKS bootstrapping core. The first diagram isolates the
key semantic invariant: coefficient split does not preserve the original
big-ring canonical slots, but after applying \(\CtwoS\), the resulting
coefficient-packed coordinates decompose by residue class. The second
diagram shows the full split-compatible sandwich.

\section{Core Split--\(\CtwoS\) Commutative Square}

The key compatibility is
\[
\left(
\bigoplus_{r=0}^{k-1}\CtwoS_n
\right)
\circ
\Split_k
=
\Pi_k
\circ
\CtwoS_N.
\]

Equivalently, the following square commutes.

\begin{figure}[h]
\centering
\resizebox{0.78\textwidth}{!}{%
$
\begin{tikzcd}[
  ampersand replacement=\&,
  column sep=huge,
  row sep=large
]
\Can_N
  \arrow[r, "\Split_k"]
  \arrow[d, "\CtwoS_N"'] \&
\displaystyle\bigoplus_{r=0}^{k-1}\Can_n
  \arrow[d, "\displaystyle\bigoplus_{r=0}^{k-1}\CtwoS_n"] \\
\Coeff_N
  \arrow[r, "\Pi_k"'] \&
\displaystyle\bigoplus_{r=0}^{k-1}\Coeff_n
\end{tikzcd}
$
}
\caption{
Core semantic commutative square. The split operation does not preserve
the original big-ring canonical slots. However, after applying
\(\CtwoS\), the resulting coefficient-packed coordinates decompose
exactly by residue class. This is the invariant used by the factorized
bootstrapping core.
}
\label{fig:split-c2s-square}
\end{figure}

The top route first performs coefficient split at the message level and
then applies leaf \(\CtwoS_n\). The bottom route first applies the
large-ring \(\CtwoS_N\) and then groups coefficient-packed coordinates by
residue class using \(\Pi_k\). The two routes produce the same leaf
coefficient-packed representation.

\section{Full Split-Compatible Sandwich}

The ideal large-ring core is
\[
\BTS^{\mathrm{core}}_N
=
\StwoC_N
\circ
\EvalMod_N
\circ
\CtwoS_N.
\]

The leaf core is
\[
\BTS^{\mathrm{core}}_n
=
\StwoC_n
\circ
\EvalMod_n
\circ
\CtwoS_n.
\]

The factorization theorem states that, under the normalized coefficient-split
semantics,
\[
\StwoC_N
\circ
\EvalMod_N
\circ
\CtwoS_N
=
\Merge_k
\circ
\left(
\bigoplus_{r=0}^{k-1}
\StwoC_n
\circ
\EvalMod_n
\circ
\CtwoS_n
\right)
\circ
\Split_k.
\]

The proof is visualized by the following diagram.

\begin{figure}[h]
\centering
\resizebox{\textwidth}{!}{%
$
\begin{tikzcd}[
  ampersand replacement=\&,
  column sep=large,
  row sep=large
]
\Can_N
  \arrow[r, "\CtwoS_N"]
  \arrow[d, "\Split_k"'] \&
\Coeff_N
  \arrow[r, "\EvalMod_N"]
  \arrow[d, "\Pi_k"'] \&
\Coeff_N
  \arrow[r, "\StwoC_N"]
  \arrow[d, "\Pi_k"'] \&
\Can_N \\
\displaystyle\bigoplus_{r=0}^{k-1}\Can_n
  \arrow[r, "\displaystyle\bigoplus_{r=0}^{k-1}\CtwoS_n"'] \&
\displaystyle\bigoplus_{r=0}^{k-1}\Coeff_n
  \arrow[r, "\displaystyle\bigoplus_{r=0}^{k-1}\EvalMod_n"'] \&
\displaystyle\bigoplus_{r=0}^{k-1}\Coeff_n
  \arrow[r, "\displaystyle\bigoplus_{r=0}^{k-1}\StwoC_n"'] \&
\displaystyle\bigoplus_{r=0}^{k-1}\Can_n
  \arrow[u, "\Merge_k"']
\end{tikzcd}
$
}
\caption{
Factorized CKKS bootstrapping core as a split-compatible sandwich.
The left square is the Split--\(\CtwoS\) compatibility, the middle square
is component-wise EvalMod compatibility, and the right rectangle is the
\(\StwoC\)--Merge compatibility. Together, these compatibilities imply
the factorization of the large-ring bootstrapping core into leaf cores.
}
\label{fig:core-factorization-ladder}
\end{figure}

The left square expresses
\[
\left(
\bigoplus_{r=0}^{k-1}\CtwoS_n
\right)
\circ
\Split_k
=
\Pi_k
\circ
\CtwoS_N.
\]

The middle square expresses the component-wise EvalMod compatibility
\[
\left(
\bigoplus_{r=0}^{k-1}\EvalMod_n
\right)
\circ
\Pi_k
=
\Pi_k
\circ
\EvalMod_N.
\]

The right rectangle expresses
\[
\Merge_k
\circ
\left(
\bigoplus_{r=0}^{k-1}\StwoC_n
\right)
\circ
\Pi_k
=
\StwoC_N.
\]

Composing these three compatibilities gives the factorized core identity.

\end{document}